\section*{Chapter~3~~Research~Problem}
\setcounter{section}{3} \setcounter{subsection}{0} 
\setcounter{table}{0} \setcounter{figure}{0} \setcounter{equation}{0}\setcounter{definition}{0}
\addcontentsline{toc}{section}{\textbf{Chapter~3~~Research~Problem}}

\noindent This Chapter is dedicated to outlining the core objectives of this project, structured around three major questions that it seeks to address. The intent of this chapter is to present these pivotal questions, providing a clear statement of the problems this project aims to solve within the domain of BLE network security. We will also delineate our project's contributions to the field, underlining the potential impact and advancements our findings offer to the existing body of knowledge.

\subsection{Problem I: Real-Time Spoofing Detection}

\noindent \textbf{$\bullet$ Problem I: How can we achieve real-time, precise analysis to efficiently identify advanced spoofing attackers?}

The ability to identify spoofing attacks as they occur is paramount in safeguarding BLE networks. Real-time analysis ensures that security measures can react instantly to potential threats, mitigating risks before they materialize into breaches. Precision in detection is equally crucial to avoid the costs associated with false alarms. By solving this problem, the positive impact would be two-fold: the creation of a more secure network environment and the assurance of user trust through the reliable protection of sensitive data.

\subsection{Problem II: Efficiency and Impact Minimization}

\noindent \textbf{$\bullet$ Problem II: How do we ensure that our detection mechanism remains efficient while minimizing its side effects on system operations and energy consumption?}

An efficient detection mechanism must operate with minimal impact on the system it protects. Excessive energy consumption or operational delays can be as detrimental as the threats they aim to prevent. Achieving this balance is crucial for the viability of security solutions in energy-constrained environments. By addressing this challenge, the resultant detection mechanism can be widely adopted, ensuring broad-scale security that is both effective and sustainable.


\subsection{Problem III: Reproducibility and Benchmarking}

\noindent \textbf{$\bullet$ Problem III: How do we ensure the reproducibility of our proposed solution and establish a new benchmark dataset to foster progression in this field of research?}

Reproducibility is the cornerstone of scientific advancement. Ensuring that our solution can be independently verified by other researchers reinforces the validity of our findings. Furthermore, by establishing a new benchmark dataset, we contribute a valuable resource that catalyzes further research. This not only demonstrates the practicality of our solution but also propels the field forward, providing a foundation for future innovations in BLE network security.

\vspace{10pt}

Each of these challenges pertains to a vital facet of BLE network security. Detailed resolutions to these challenges will be elucidated in the subsequent chapter.

\newpage